\documentclass[11pt]{article}
\usepackage[utf8]{inputenc}
\usepackage[T1]{fontenc}
\usepackage{lmodern}
\usepackage[margin=1in]{geometry}
\usepackage{csquotes}
\usepackage[style=authoryear,backend=biber,sorting=nyt,maxbibnames=99,giveninits=true,uniquename=init]{biblatex}
\usepackage[hidelinks]{hyperref}
\addbibresource{references.bib}
\begin{document}
\section*{Numerical Modeling of Pipe Flow Dynamics}

\subsection{Overview}

Numerical and analytical modeling of pipe flow dynamics in subsurface (tile) drainage
spans a continuum of scales and physical processes, from the movement of water toward
and into individual perforated pipes, through the hydraulics of flow within the conduit
network, to the measurement of discharge at control structures and outlets. Across this
theme, the dominant methodological signature is validation: models are consistently
tested against sand-tank and laboratory experiments, analytical benchmarks, field
observations, and independent numerical solutions to confirm accuracy in predicting
inflow/outflow dynamics, water-table behavior, and water mass balance. Several
consolidated lines of evidence emerge. First, the classical governing equations---the
Richards equation and its Boussinesq-based approximations---reproduce observed
water-table drawdown and cumulative drainage/subirrigation volumes when a nonconstant
drainable porosity is accounted for \parencite{TANG1977,Xin2016c}. Second, analytical
drain-spacing and single-tile flow solutions closely track numerical simulations and
observed heads, and newer formulations reduce reliance on the difficult-to-estimate
drainable porosity parameter \parencite{Cooke2001,KUMAR1994,SINGH1992,FathAllah2025,Chahar2010}.
Third, entrance resistance and effective radius---the parameters governing how readily
water crosses the pipe wall---have been quantified for a range of perforation geometries
and envelopes, with sock-wrapped pipes approaching the behavior of a fully open pipe
\parencite{Ghane2022,Ghane2022b,Gaj2020,Oyarce2017}. Fourth, within-conduit hydraulics,
including spatially varied flow, friction resistance, and outlet effects, have been
modeled and benchmarked against experiments with sub-5\% error in leading cases
\parencite{Bayat2011,Kouchakzadeh2011,Abdurrasheed2019,Graber2004,Nédélec2008}. Finally,
robust head--discharge relationships for V-notch and compound weirs, along with
vision-based flow detection, provide the measurement backbone for validating drainage
models and quantifying conservation-practice performance
\parencite{Christianson2019,BinShokrana2021,Katuwal2024,Napierala2023,Popek2021,Young2023}.
This section synthesizes these established results, emphasizing quantified outcomes and
the conditions under which each modeling approach has been confirmed.

\subsection{Governing Equations for Water-Table and Drainage Dynamics}

\subsubsection{Richards- and Boussinesq-Based Solutions}

The most physically complete descriptions of drainage and subirrigation solve the
Richards equation, which fully represents saturated and unsaturated water movement during
both processes \parencite{TANG1977}. In sand-tank validation, predicted water-table
drawdown from the Richards equation was somewhat slower than observed but in generally
good agreement; the profile drained slightly more rapidly than predicted, whereas for
subirrigation the predicted--observed agreement was better than for drainage
\parencite{TANG1977}. Predicted and observed inflow/outflow rates and cumulative flow
volumes agreed closely, with the predicted rate slightly lower than observed for drainage
and slightly higher than observed for subirrigation \parencite{TANG1977}. Critically,
approximate Boussinesq-based solutions were in close agreement with the Richards-equation
solutions and with experiment \emph{when a variable (nonconstant) drainable porosity was
considered} \parencite{TANG1977}. The overall conclusion is that the Richards equation is
more accurate, but Boussinesq approximations, corrected for nonconstant drainable
porosity, remain useful for design when reliable saturated hydraulic conductivity and
drainable porosity are available; by suitable modification the computer solution can also
accommodate irregular boundary/initial conditions and soil stratification
\parencite{TANG1977}.

The role of the unsaturated (vadose) zone in setting the effective drainable porosity is
developed further in the Boussinesq-equation (BE) analyses of \textcite{Xin2016c}. Using
an effective specific yield formulation to mimic the vadose-zone effect, the authors
report $S_y = 0.0512$ for cases 1 and 3 and $0.1057$ for cases 2 and 4, and derive a new
analytical solution separately for the BE model to remove large errors that arise when the
unsaturated zone is neglected \parencite{Xin2016c}. The two BE-model solutions predicted
almost identical non-dimensional cumulative water drainage over time across different
$l_1/l$ values, indicating that the new analytical solution is reasonably accurate across
scales, and for the four experimental runs the adjusted-$S_y$ BE model, the analytical
solution, and the SUTRA numerical model predicted similar total water drainage at the end
of the process---good agreement in water mass balance \parencite{Xin2016c}. The BE model
gave reasonably good agreement with the experimental data for cases 2 and 4. A recognized
limitation is that the BE model with an effective specific yield predicted a quicker
groundwater-table decline and could not capture the dynamic drainable-porosity effect
caused by vadose-zone development; for case 4 the analytical solution predicted slightly
quicker drainage than the numerical BE model, attributed to approximations (notably a
Taylor expansion of the initial boundary condition) \parencite{Xin2016c}. Overall, the
Boussinesq equation with an effective specific yield provides a reasonably good
description of the drainage process, particularly for total water drainage
\parencite{Xin2016c}.

\subsubsection{Preferential Flow and Constitutive Relationships}

At the field scale, the fidelity of tile-flow prediction depends strongly on the soil
hydraulic functions used. \textcite{Mohanty1997} introduced new piecewise-continuous,
bimodal soil-water-retention and hydraulic-conductivity functions, which predicted
preferential flow to the tile drain reasonably well following all irrigation events and
performed better than the unimodal van Genuchten--Mualem functions in describing the
observed flow regime. Predicted flow using the bimodal functions matched observed flow
rates reasonably well and far better than the unimodal case, minimized over- and
under-predictions across the different flow domains, and better captured the timing and
magnitude of the preferential-flow peaks \parencite{Mohanty1997}. This body of work
established that capturing preferential (macropore) flow in tile drains requires
constitutive relationships beyond a single unimodal curve.

The recession behavior of the water table following recharge was addressed analytically
by \textcite{VANDELEUR1959}, whose mathematical approach accorded with field-experiment
conclusions and was successfully checked by a sandbox experiment (in which radial
resistance was absent). Tail recession was described as beginning only after a period
following the cessation of vertical percolation into the saturated zone; before that, the
recession curve depends on the rates and durations of preceding replenishment
\parencite{VANDELEUR1959}. Radial resistance near the drain could appear as a modified
value of the equivalent depth $D$, suggested to be high at one site and almost absent at
another \parencite{VANDELEUR1959}.

\subsection{Analytical Drain-Spacing and Single-Tile Flow Solutions}

\subsubsection{Single-Tile and Region-of-Influence Formulations}

For flow to a single tile, \textcite{Cooke2001} proposed analytical solutions that closely
agreed with numerical (finite-element) simulations and were judged adequate for practical
description of single-tile flow. The key conceptual advance is replacing the usual drain
spacing with a variable \emph{region of influence} $L$---corresponding to points in the
finite-element domain where saturation was at least 98.5\%---which makes the equations
applicable to irregular or randomly spaced tile systems and eliminates the need to specify
a spacing parameter \parencite{Cooke2001}. Without accretion, analytical solutions using
drainable porosity derived from water-retention curves gave root-mean-squared deviations
of $0.085$ cm$^2$ per day for silty clay loam and $0.060$ cm$^2$ per day for clay;
deviations between analytical and numerical flow were greatest near the start of
simulations, when the initially flat water table was being drawn down to the tile
\parencite{Cooke2001}. The resulting equations are simple, universally applicable, and can
be incorporated into existing models without appreciable added complexity or computing
cost \parencite{Cooke2001}.

\subsubsection{Transient Water-Table and Parameter-Estimation Methods}

\textcite{SINGH1992} modified the steady-state drainage solutions of Ernst, Dagan, and van
Beers to predict the transient water table midway between drains, integrating them
alongside the Hooghoudt, Toksöz--Kirkham, and van Schilfgaarde equations. At the Mundlana,
India site these integrated equations predicted hydraulic heads reasonably close to
observed values for 50, 67, and 84 m drain spacings and agreed well with the transient
solutions of van Schilfgaarde and of Bouwer and van Schilfgaarde \parencite{SINGH1992}.
For the Nwa and Twocock data, all equations predicted larger spacings than the actual
40.23 m while maintaining the same order of performance across both experimental periods;
for the El-Mowelhi and Hermsmeier data, the predicted time for a 0.8 m water-table drop
was two to three times the actual time in the first year and four to six times in the
second year, with some equations agreeing more closely than others
\parencite{SINGH1992}.

Complementary parameter-estimation work by \textcite{KUMAR1994} in India proposed
analytical solutions that agreed closely with the Glover--Dumm solution, with reported
errors below 10\% for periods up to 15 days. A computer-simulation technique produced high
agreement between observed and predicted dimensionless hydraulic heads (correlation
coefficients of 0.9449 to 0.9931), yielded $KD/f$ estimates close to graphical-technique
means, and was concluded to be ``by far and most appropriate'' for evaluating the
transmissivity-to-drainable-pore-space ratio \parencite{KUMAR1994}. Groundwater-table
predictions using inverse-technique parameters matched observations better than those
using field-investigation parameters, and the graphical and computer-simulation techniques
were the preferred approaches for predicting $KD/f$ \parencite{KUMAR1994}. A
discharge-based inverse technique had limited applicability---the Mundlana data were not
amenable to it and the trial-and-error procedure was time consuming unless discharge
measurements were widely separated in time---and the proposed drainable-pore-space
equation generally gave lower $f$ values \parencite{KUMAR1994}. The work also proposed
practical procedures to separate transmissivity and drainable pore space once their ratio
is known, filling a gap in the available theory \parencite{KUMAR1994}.

\subsubsection{Design Equations Reducing Reliance on Drainable Porosity}

A recurring practical obstacle in unsteady-state design is the drainable porosity of soil
($f$). \textcite{FathAllah2025} addressed this with the Safaa equation, which avoids using
$f$---described as the greatest obstacle in unsteady-state design---and replaces it with a
new constant $C_s$, while adding $C_{sh}$ to represent the shape of the initial water
table. Validation with Haress pilot-area data (Nile Delta, Egypt) showed that the lateral
spacing from the Safaa equation was closer to the actual spacing than that from Glover's
equation, and the exponential curve plotted with the Safaa equation fell within the range
of the actual curves \parencite{FathAllah2025}. The Safaa spacing is greater than
Glover's, making the design more economic, and the study concludes the equation provides
more accurate and economic results for subsurface drainage design in the Delta region
\parencite{FathAllah2025}.

For minimum-cost ditch-drainage design, \textcite{Chahar2010} used nonlinear data fitting
to derive an explicit equation for optimal ditch spacing, avoiding numerical computation
of improper integrals, simultaneous solution of nonlinear equations, and formal
optimization. The fitted equation had errors within 7\% over the practical parameter
ranges: for Case 1 it gave 4.40 m versus 4.586 m from the optimization program (4.1\%
error), and for Case 2, 12.035 m versus 12.514 m (3.83\% error) \parencite{Chahar2010}.
Optimal spacing increases with the earthwork-cost parameter and decreases with the
pumping-cost parameter, and the explicit forms give close approximations across the
practical ranges considered \parencite{Chahar2010}.

\subsection{Entrance Resistance, Effective Radius, and Perforation Design}

\subsubsection{Effective Radius of Perforated and Enveloped Pipes}

The ease with which water crosses the pipe wall is characterized by the effective radius
and the entrance resistance factor. \textcite{Ghane2022b} quantified effective radii for
Michigan pipes: 3- or 4-row regular-perforated pipes ranged from 0.3 to 0.9 cm (average
0.6 cm) and 8-row regular-perforated pipes reached 1.9 cm, while 4-row sand-slot pipes
ranged from 0.1 to 0.7 cm (average 0.4 cm) and the 8-row sand-slot pipe reached 1.6 cm.
Wrapping regular-perforated pipes with a knitted-sock envelope increased the effective
radius dramatically to 5.7--6.0 cm (average 5.9 cm) and reduced entrance head loss to near
zero \parencite{Ghane2022b}. The most efficient way to increase effective radius was
increasing the number of longitudinal rows of perforations, followed by increasing
perforation length, whereas increasing perforation width had minimal effect
\parencite{Ghane2022b}. These hydraulic gains carry design and cost implications:
sock-wrapped pipes allowed wider optimum drain spacing than sand-slot pipes, but 8-row
sand-slot pipes had a lower initial system cost than sock-wrapped pipes across all factor
combinations at the same drainage rate and depth; and for soils without a sedimentation
problem, 8-row regular-perforated pipes allowed wider spacing and lower initial cost than
4-row pipes \parencite{Ghane2022b}.

A companion study \parencite{Ghane2022} experimentally measured the effective radius of a
commonly manufactured sock-wrapped pipe for the first time and found close agreement
between measured and theoretical effective radii for the sock-wrapped and sand-slot pipes,
validating the Dierickx (1999) theoretical method. The sock-wrapped pipe had much lower
entrance head loss than the 8-row and 4-row sand-slot pipes, giving it the largest
effective radius and a higher estimated drain inflow than both sand-slot pipes; its very
small entrance resistance factor meant it functioned almost as a completely open pipe
without walls \parencite{Ghane2022}. The knitted-sock envelope met the hydraulic,
anti-clogging, and retention criteria for the sand used, and under free drainage with
ponding, water entered the sock-wrapped pipe only through perforations in the bottom half
of the pipe \parencite{Ghane2022}.

\subsubsection{Entrance Resistance in Corrugated Versus Plain Pipes}

Using a numerical model calibrated with sand-tank experiments, \textcite{Gaj2020}
estimated entrance resistance ($\alpha_e$) and delivery ratio ($Q/Q_0$) for perforated
pipes and developed six nonlinear multiple-regression models to predict these from
perforation design variables---the first reported computation of entrance resistance for
corrugated pipes with circular holes. For corrugated pipes with circular holes, $\alpha_e$
was higher than for plain-wall pipes with the same perforation configuration because
streamlines converge in the corrugation valley; for the case $N=4$ and
$A_p = 60$ cm$^2$/m, $\alpha_e$ was 23\% higher and $Q/Q_0$ 20\% lower in corrugated than
in plain-wall pipes \parencite{Gaj2020}. For the same opening area and configuration,
$\alpha_e$ was 51\% higher for circular holes than for rectangular slots, while $Q/Q_0$
was marginally greater for circular holes; the drawdown at the midpoint between two
parallel drains did not differ significantly between holes and slots for $N=4$ and
$A_p = 60$ cm$^2$/m \parencite{Gaj2020}. The consolidated conclusion is that rectangular
slots are more hydraulically advantageous than circular holes in the valleys of corrugated
pipes \parencite{Gaj2020}.

\textcite{Oyarce2017} gave experimental validity to Dierickx's mathematical model through
discharge measurements in a physical prototype and developed a calculation tool to
evaluate water-entrance resistance across pipe radii, perforation densities, and screen
presence. Water extraction by a perforated pipe was proportional to extraction area, but
increasing perforation density reduced pipe support and mechanical-impact resistance
\parencite{Oyarce2017}. The entrance-resistance factor depended on perforation density and
arrangement and on the presence of an enveloping screen; a screen reduced the resistance
factor and increased the effective radius, whereas identical pipes without an envelope did
not drain the substrates because the perforations clogged almost instantly
\parencite{Oyarce2017}. For small pressure heads, discharge increased with perforation
density, but for larger heads the effect was less relevant; the optimal distance between
adjacent parallel drainpipes was smaller when entrance resistance was considered,
especially for low saturated hydraulic heads \parencite{Oyarce2017}. Overall, smaller
diameter pipes with higher perforation density and an enveloping screen provided greater
discharge capacity than larger-diameter pipes---subject to a balance between mechanical
strength and discharge capacity---identifying an economic/design advantage for smaller
pipes with few, small-width perforations combined with a screen envelope
\parencite{Oyarce2017}.

\subsubsection{Structural Response of Perforations to Traffic Loads}

The mechanical counterpart to perforation hydraulics is the structural response of
corrugated HDPE pipes to field traffic, modeled by \textcite{Gaj2021}. Vertical stress
transmission from wheel loads was largest for the loaded grain wagon and smallest for the
tractor; at depths greater than 1.6 m, $\sigma_{vi}$ was below 35 kPa for the three
wheel-load cases \parencite{Gaj2021}. A bonded interface-friction condition generally
induced larger wall stresses than a smooth condition, though not statistically
significantly across all soil types and load cases; for non-perforated pipes the loaded
(LL) cases produced greater stresses than the empty/half-loaded cases across all soil
types, with the critical case being a loaded grain wagon on silt loam \parencite{Gaj2021}.
Failure by ductile yielding of non-perforated corrugated HDPE pipes was highly unlikely
under the investigated field-traffic, soil, and water-table conditions, and vertical
deflection remained within the 5\% performance limit under all loading, although
perforations can raise the risk of ductile-yield failure before deflection becomes
limiting \parencite{Gaj2021}. Circular-hole perforations substantially increased local
wall stresses---5.6 mm-diameter holes exceeded the yield strength at burial depths
shallower than 0.9 m under the critical load case in silt loam---whereas rectangular slots
produced lower stress concentrations when the perforation width $w_{perf} \le 0.5 w_v$,
with slot length having little effect compared with slot width \parencite{Gaj2021}. Axial
perforation spacing did not materially affect stress distribution, so perforations may be
placed in every corrugation valley without direct axial-stress interaction; the simulated
stresses were used to establish perforation-dimension relationships that allow perforations
to be incorporated into the structural design of buried corrugated pipes
\parencite{Gaj2021}. Consistent with the hydraulic findings, rectangular slots also had a
lower ductile-yielding failure risk than circular holes provided $w_{perf} \le 0.5 w_v$
\parencite{Gaj2021}.

\subsubsection{Discharge to Improved Drains Under Ponding}

\textcite{Tao2019} proposed a general discharge formula for improved subsurface drainage
with ponded water in homogeneous, isotropic soil, verified by laboratory experiment, field
test, and numerical simulation. Under two boundary assumptions, Method 1 outperformed
Method 2 in the numerical simulation, with calculated values closer to the 1:1 line; the
critical boundary between saturated and unsaturated filter conditions was
$(b_0 + h_0)/d = 0.6$ for Method 1 and $h_0/d = 0.4$ for Method 2 \parencite{Tao2019}. For
formula simplification, a drain-spacing-to-impermeable-layer-depth ratio of 3 was the
critical condition separating the two simplified equations, and when the impermeable-layer
depth is unknown, a fallback equation can roughly estimate discharge \parencite{Tao2019}.

\subsection{Within-Conduit Hydraulics and Network Flow}

\subsubsection{Spatially Varied Flow and Discharge Capacity}

Flow within drainage conduits with lateral inflow is a spatially varied flow (SVF)
problem. \textcite{Bayat2011} calibrated an SVF model by an inverse method, proposing an
average Manning coefficient of 0.02 for the tested corrugated pipes and local head-loss
coefficients $K_{ud} = 1.0$ and $K_{bd} = 1.0$, supported by comparison between computed
and observed water-surface elevations. Ignoring local head loss in the manhole computation
produced underestimation with a maximum error of $-20\%$ \parencite{Bayat2011}. For the
example network, the SVF model gave a total discharge capacity of 6.41 l\,s$^{-1}$ when
$y/D$ was close to 0.94 in the main pipe, and 5.12 l\,s$^{-1}$ when no segment was allowed
to pressurize; using modified traditional relationships for the same network gave much
larger capacities of 12 and 20.81 l\,s$^{-1}$---roughly two to four times the SVF
value---indicating a large discrepancy relative to the physically based SVF method
\parencite{Bayat2011}.

The generalized capacity problem was treated by \textcite{Graber2004}, who presented a
generalized numerical solution for circular conduits with free discharge and uniform
inflow over a wide range of parameters, and a companion solution for a rectangular
horizontal channel judged of particular practical value. Manning's equation alone is not
generally adequate for subsurface-drain capacity, and Van der Beken's method is likewise
inadequate; computer-based numerical methods are recommended more generally
\parencite{Graber2004}. In the worked example, a 0.102 m pipe provided insufficient
capacity while a 0.152 m pipe met the required design capacity, and a multi-diameter
interceptor-drain design (14-, 12-, 10-, 8-, and 6-in. segments) was produced by computer
calculation \parencite{Graber2004}. Supercritical flow should usually be avoided in
subsurface drains because it uses pipe capacity uneconomically, produces relatively high
velocities, and increases turbulence through hydraulic jumps \parencite{Graber2004}.

\subsubsection{Flow Resistance in Corrugated Pipes}

Resistance to flow depends strongly on flow regime. \textcite{Kouchakzadeh2011} compared
pressurized flow, gradually varied flow (GVF), and SVF in corrugated pipes. The pressurized
setup revealed Darcy--Weisbach friction factors of 0.065 and 0.04 for the 0.115 m and
0.182 m diameter pipes, with corresponding Moody-diagram relative roughness of 0.04 and
0.012 \parencite{Kouchakzadeh2011}. Reported Manning's coefficients were 0.025 and 0.020
for GVF and 0.035 and 0.030 for SVF, for the two diameters respectively; the GVF resistance
coefficient exceeded that of the pressurized regime, and SVF increased resistance by
40--50\% relative to GVF \parencite{Kouchakzadeh2011}. These results provide Manning's $n$
values for corrugated pipes used in sewer and drainage networks \parencite{Kouchakzadeh2011}.

\subsubsection{Outlet Effects and Flow-Control Modules}

Flow behavior at conduit outlets was studied by \textcite{Nédélec2008}, extending
knowledge of free-pipe outlet flow. The conduit outlet may cause non-negligible tail-water
effects upstream when high pipe flow combines with high channel discharge capacity, and
this tail-water effect may increase with flow rates when pipe and channel flows are
concomitant and proportional \parencite{Nédélec2008}. Combined pipe slope and downstream
tail-water effects may cause a transition from free-surface to pressurized flow, and the
system can exhibit two forms of transitional cavity flow with possibly different transient
behaviors \parencite{Nédélec2008}. The work identified and illustrated ten flow patterns
in the channel and pipe; the outlet energy-loss behavior resembled junctions of two pipes
of different diameters and compared rather well with a more general junction model,
suggesting possible generalization subject to further study of friction losses
\parencite{Nédélec2008}.

The effect of an in-channel flow-control module was modeled in three dimensions using
FLOW-3D by \textcite{Abdurrasheed2019}. The numerical results closely matched laboratory
experiments in both free flow without module (FFWM) and flow with module (FWM), with
percentage error below 5\%, and the experimental results strongly agreed with the numerical
values \parencite{Abdurrasheed2019}. The REDAC module altered the flow pattern and created
resistance, increasing the downstream depth-averaged velocity from 43.50 cm/s to about
46.50 cm/s and decreasing water depth from 7.75 cm to 6.50 cm, while raising Manning's
roughness upstream relative to FFWM \parencite{Abdurrasheed2019}. Mass conservation was
maintained, with inflow and outflow both coinciding at 4750 cm$^3$/s (0.00475 m$^3$/s) in
both cases \parencite{Abdurrasheed2019}. The FLOW-3D approach modeled the module behavior
with greater certainty, visualized flows in 2D and 3D, and can characterize larger module
arrangements and different flow conditions without expensive laboratory experiments
\parencite{Abdurrasheed2019}.

\subsubsection{Network Routing Through Risers and Berms}

At the watershed scale, \textcite{Gupta2019} presented a new toolbox (CoBAGNPS, based on
the AGNPS model) for simulating water and sediment control basins (WASCoBs), which
successfully reproduced the required outputs and routed flow through WASCoBs and the
drainage network. For a single WASCoB, there was no substantial difference in flow
characteristics between HB and PI pipe risers, while a third riser type (HBS) produced a
shorter peak discharge and drained slightly more slowly; at the DFTILE outlet the three
riser types produced different peak discharges for two storm events but with overall
reasonable routing behavior \parencite{Gupta2019}. For extreme design storms, the installed
150 mm tile drain was not efficient enough; larger pipes improved drainage, and a 375 mm
pipe was recommended for 2-, 5-, and 10-year 24-h storm events, with 200 mm and 375 mm
diameters significantly enhancing efficiency and reducing ponding time behind the berms
\parencite{Gupta2019}. Direct routing from the DFTILE outlet to the watershed outlet had
only a slight effect on peak flow but increased the recession limb of the hydrograph
\parencite{Gupta2019}.

\subsection{Flow Measurement at Control Structures and Outlets}

\subsubsection{V-Notch Weir Calibrations}

Accurate discharge measurement underpins the validation of all drainage models and the
assessment of conservation practices \parencite{Christianson2019,Katuwal2024}. For a new
stainless-steel-edged 45\textdegree{} V-notch weir in AgriDrain inline water-level control
structures, \textcite{Christianson2019} found no significant differences in fitted weir
parameters across the two structure sizes (15 cm and 25 cm) or the three weir placements.
The fitted equation was $Q = 0.011 H^{2.28}$ ($Q$ in l\,s$^{-1}$, $H$ in cm), equivalently
$Q = 1.44 H^{2.28}$ in U.S. customary units, with a combined-model coefficient of
determination of 0.992 and a standard error of estimate of 0.135
\parencite{Christianson2019}. The standard deviation of measured flow for each weir-height
and structure-size combination ranged from 0.1\% to 2.9\% of the mean, and the equation was
valid only up to a maximum water depth of 16.5 cm, above which a compound weir equation
should be used \parencite{Christianson2019}.

\textcite{BinShokrana2021} developed an empirical flow equation for a metal-edge
sharp-crest V-notch weir using a top-down approach that gave a V-notch apex height of
31.7 cm (from $b = 149.8$ cm and $c = 118.1$ cm). A combination of the weighing method and
a flow meter produced the stage--discharge relation, and least-squares regression showed a
strong head--flow relationship ($R^2 = 0.999$) with the empirical equation
$Q = 0.749 H^{2.25}$, calibrated for flow rates up to 409.42 L/min at a head of 16.3 cm
\parencite{BinShokrana2021}. This equation produced consistently higher flow estimates than
the \textcite{Christianson2019} equation, whereas the USGS equation underestimated flow
across all head measurements \parencite{BinShokrana2021}.

For a PC-edge sharp-crest 22.5\textdegree{} V-notch weir, \textcite{Napierala2023} produced
empirical stage--discharge equations for downstream-face (chamfer) slopes of 30\textdegree,
45\textdegree, and 60\textdegree. The 45\textdegree{} chamfer gave the strongest fit and
lowest error ($R^2 = 0.9981$, RMSE $= 0.255$ dm$^3$\,min$^{-1}$), with the 30\textdegree{}
and 60\textdegree{} equations also fitting well ($R^2 = 0.9955$ and $0.9980$); residuals
were generally within $\pm 4\%$, and among previously reported equations White's fit the
data best for all chamfered angles \parencite{Napierala2023}. These equations can be used
for any cylindrical control structure provided flow passes through the V-notch without
submergence and can support nutrient-load-reduction estimates in subsurface drainage
studies, with the 45\textdegree{} chamfer the best choice \parencite{Napierala2023}.

\subsubsection{Compound Weirs for Overtopping Flows}

Because V-notch equations lose validity once flow overtops the notch, compound weir
formulations are required for high flows. \textcite{Katuwal2024} developed a compound weir
equation ($Q_{CW}$) that estimated overflow discharge accurately, whereas the existing
combined V-notch-plus-rectangular method ($Q_{V+R}$) consistently underestimated it. Across
three Agri Drain control-structure sizes, $Q_{CW}$ produced MAPE below 5\% and PBIAS
between $-5.4\%$ and $0.7\%$, while $Q_{V+R}$ had PBIAS above 15\% and MAPE above 13\%
\parencite{Katuwal2024}. The V-notch calibration relationships varied with structure size,
confirming that head--discharge relationships are affected by structure size; against
published USBR compound-weir data, $Q_{CW}$ gave estimates comparable to measured values
while $Q_{V+R}$ remained consistently lower \parencite{Katuwal2024}. Importantly, using
$Q_{CW}$ instead of $Q_{V+R}$ increased the estimated annual discharge and nitrate load
diverted to a saturated buffer, demonstrating that overflow underestimation can materially
bias nutrient-load estimates \parencite{Katuwal2024}.

\subsubsection{Adjustable Damming Regulators}

\textcite{Popek2021} calibrated two water-damming regulators for drains. The funnel
regulator had an effective flow of 0.17--0.25 l\,s$^{-1}$ over $H_f = 47$--71 cm, with a
linear relation between effective flow and overflow-edge height ($R^2 = 0.960$); when inflow
increased up to 5.5 l\,s$^{-1}$, the water layer above the overflow edge rose only
0.12--3.0 cm, allowing a maximum damming adjustment of up to $+3.0$ cm relative to the set
height \parencite{Popek2021}. For the hole regulator, the effective flow required for
damming was 0.009 l\,s$^{-1}$ at $H_h = 20$ cm, 0.012 l\,s$^{-1}$ at 40 cm, and
0.015 l\,s$^{-1}$ at 60 cm; flows up to 0.25 l\,s$^{-1}$ produced similar discharges among
holes and water-table elevations within about 11--15 cm above the hole edge, above which
damming height became more variable, and the reduced-flow curve was developed for
intensities up to 0.65 l\,s$^{-1}$ \parencite{Popek2021}. The funnel regulator suits
holding water at flows at or above the effective flow and ensures a stable damming ordinate
across a wide flow range, while the hole regulator suits low flow rates up to
0.25 l\,s$^{-1}$; both have uncomplicated structure and adapt to existing drainage devices
\parencite{Popek2021}.

\subsubsection{Vision-Based Outlet Monitoring}

Addressing a gap in outlet-monitoring technology, \textcite{Young2023} developed a
low-cost, vision-based method to characterize outlet flow events by occurrence, duration,
and rate from free-standing outlet points without in-situ instrumentation. The approach
correctly detected outlet flow in low, medium, and high flow test cases and did not
misclassify hand or object motions as flow when none was present; occurrence detection was
delayed relative to ground truth by an average of 1.0 s and flow duration was underestimated
by 2.23\% on average \parencite{Young2023}. Steady-state flow-rate estimates were lower than
ground truth for all three conditions, with errors of 5.9\%, 9.2\%, and 11.2\%, so flow
rates were estimated within 12\% or less of ground truth under laboratory conditions
\parencite{Young2023}. The study validated the feasibility of visually monitoring tile
outlet flow and released an open-source application to extract flow-rate data from recorded
video \parencite{Young2023}.

\subsection{Coupling Drainage Hydraulics to Buffer Performance}

Numerical modeling of pipe flow also feeds into the design of edge-of-field practices that
receive drainage water. \textcite{McEachran2020} balanced the time to remove nitrate
against the time for water to flow through a saturated riparian buffer (SRB) to derive a
quantitative expression for optimal buffer width. For all six study sites the optimal width
was smaller than the current SRB width; widths at four of the six sites exceeded the upper
limits of the optimal ranges, implying that nitrate-removal effectiveness could be about 50
to 120\% greater at those sites, while widths at two sites fell within the optimal range
with near-optimal effectiveness \parencite{McEachran2020}. Two sites had optimal widths
below the current minimum buffer width in USDA-NRCS Code 604 \parencite{McEachran2020}. The
overall uncertainty of the optimal width was large---combining all parameter uncertainties
yielded 57\% uncertainty in $L_x^*$---with the removal coefficient $\lambda$ contributing
most, followed by $K$, while $n_e$ and $\Delta h$ contributed only small amounts
\parencite{McEachran2020}. The analysis rested on first-order nitrate removal, quasi-steady-
state concentration profiles, and omission of dispersion; incorporating a site-conditioned
width into design standards would improve water quality locally and regionally
\parencite{McEachran2020}.

\printbibliography
\end{document}
